\chapter{Combinatoria de conjuntos, relaciones y funciones}

\section{Cardinal de conjuntos y cantidad de relaciones}

%\subsection{Principios de cardinal}
\begin{teoria}[title=Definion de cardinal de un conjunto]
    \textbf{sea A un conjunto, se llama cardinal de A a la cantidad de elementos distintos que tiene A, y se anota \#A. }
\end{teoria}
\begin{ejemplo}
    \# $\emptyset$ = 0 , \#\{a, b, c\}= 3  
\end{ejemplo}

\begin{teoria}[title=Propiedades de los cardinales]
    \textbf{sea A y B dos conjuntos} 
    \begin{listapropiedades}
  \item[1] Sea A un Conjuntos finitos y sea $B \subseteq A $ entonces $\#B \leq \#A $
  \item[2] $(A\cup B)$ = $\#A + \#B - \#(A \cup B)$ 
  \item[3] $\#\mu < \infty \rightarrow \#A^c = \#\mu - \#A$ 
  \item[4] $\#(A-B) = \#A - \#(A\cap B) $
  \item[5] $\# (A\Delta B) = \#A - \#B -2\#(A\cap B)$ 
  \item[6] $\# (A X B) = \#A * \#B $
  \item[7] $P(A) = 2^{\#A}$
    \end{listapropiedades}
\end{teoria}

%\subsection{Cardinales y Funciones}
\begin{teoria}[title= Cantidad de Funcion]
    \textbf{Sean A y B conjuntos finitos}
        \begin{listapropiedades}
  \item[1] $\#(f: A\rightarrow B)$, f es funcion. Entonces $\#B^{\#A}$
  \item[2] $\#(f: A\rightarrow B)$ f es una funcion Biyectiva. Entonces $n!$
  \item[3] $\#(f: A\rightarrow B)$ f es una funcion Inyectiva. Entonces $\frac{k!}{(k-n)!}$
    \end{listapropiedades}
\end{teoria}


\begin{teoria}[title= Cardinales y Funciones]
    \textbf{Sean A y B conjuntos finitos}
        \begin{listapropiedades}
  \item[1] Sea $f: A\rightarrow B$ es una funcion Inyectiva. Entonces $\#A \leq \#B$
  \item[2] Sea $f: A\rightarrow B$ es una funcion Sobreyectiva. Entonces $\#A \geq \#B$
  \item[3] Sea $f: A\rightarrow B$ es una funcion Biyectiva. Entonces $\#A = \#B$
    \end{listapropiedades}
\end{teoria}

\section{El factorial}

\section{El Numero combinatorio}

\begin{proposicion}[title=Teorema: Número Combinatorio]
    Sea $n \in \mathbb{N}_0$ y sea $A_n$ un conjunto con $n$ elementos. Para $0 \leq k \leq n$, se nota con el símbolo $\binom{n}{k}$, que se llama \textbf{número combinatorio}, y representa la cantidad de subconjuntos con $k$ elementos de un conjunto de $n$ elementos, sin importar el orden interno de cada subconjunto, y es igual a:
    \[
    \binom{n}{k} = \frac{n!}{k!(n-k)!}
    \]
\end{proposicion}

\begin{demostracion}
    Para formar una lista ordenada de $k$ elementos extraídos de un conjunto de $n$ elementos, tenemos $n$ opciones para el primer elemento, $n-1$ para el segundo, y así sucesivamente, obteniendo un total de:
    \[
    n(n-1)(n-2)\cdots(n-k+1) = \frac{n!}{(n-k)!}
    \]
    formas ordenadas (variaciones simples). Como el orden interno de esos $k$ elementos no nos importa, y cada subconjunto de $k$ elementos puede ordenarse de $k!$ maneras distintas, dividimos el total anterior por $k!$ para eliminar las repeticiones debidas al orden. De este modo, la cantidad de subconjuntos sin importar el orden es:
    \[
    \binom{n}{k} = \frac{n!}{k!(n-k)!}
    \]
\end{demostracion}

\begin{teoria}[title= Teorema del Binomio]

\end{teoria}

\begin{teoria}[title= Triangulo De Pascal]

\end{teoria}

\begin{demostracion}
    Partimos de la expresión inicial aplicando la definición de número combinatorio:
    \[
    \binom{n-1}{k-1} + \binom{n-1}{k} = \frac{(n-1)!}{(k-1)!\,[(n-1)-(k-1)]!} + \frac{(n-1)!}{k!\,[(n-1)-k]!}
    \]
    Simplificamos los factoriales en los denominadores:
    \[
    = \frac{(n-1)!}{(k-1)!(n-k)!} + \frac{(n-1)!}{k!(n-1-k)!}
    \]
    Para sumar ambas fracciones, buscamos un denominador común, que es $k!(n-k)!$. Para ello, multiplicamos y dividimos el primer término por $k$ y el segundo por $(n-k)$:
    \[
    = \frac{(n-1)! \cdot k}{k!(n-k)!} + \frac{(n-1)! \cdot (n-k)}{k!(n-k)!}
    \]
    Agrupamos todo bajo un mismo denominador y sacamos factor común $(n-1)!$ en el numerador:
    \[
    = \frac{(n-1)! \cdot [k + (n - k)]}{k!(n-k)!}
    \]
    Operamos dentro del corchete ($k - k = 0$):
    \[
    = \frac{(n-1)! \cdot n}{k!(n-k)!}
    \]
    Como $(n-1)! \cdot n = n!$, obtenemos el resultado final:
    \[
    = \frac{n!}{k!(n-k)!} = \binom{n}{k}
    \]
\end{demostracion}
